\chapter*{不定积分每日一题}

\refer{《积分的方法与技巧》金玉明\ 顾新身\ 毛瑞庭\ /\ 编著}

\exercise{求积分$\int{\sin^8{x}\d x}$}%在花括号内输入题目

\teacher{张晓鹏}%在花括号内输入你的名字

\postdate{2022年12月21日}%在花括号内输入你想要发布那一天的日期

\result{\frac{1}{2^7}\qty(\frac{\sin{8x}}{8}-8\frac{\sin{6x}}{6}+28\frac{\sin{4x}}{4}-56\frac{\sin{2x}}{2}+35x)+C}%在花括号内输入本题结果

\analysis{
    今天我们主要介绍的是三角函数积分中的倍角法.对于诸如 的三
    角函数积分，我们通过倍角法可以将高次幂的三角函数转化为一次幂的倍角三角函数.
    
    在介绍倍角法之前，我们得先介绍下今天的主要工具，棣莫弗公式：
    \formula{\qty(\cos{x}+\ii\sin{x})^n=\cos{nx}+\ii\sin{nx}}
    如果不熟悉这个的话，我们也可以用欧拉公式，这个相对熟一点.
    
    通过欧拉公式我们可以知：
    \formula{\ee^{\ii x}=\cos{x}+\ii\sin{x}}
    \formula{\ee^{-\ii x}=\cos{x}-\ii\sin{x}}

    我们令$\ee^{\ii x}=\cos{x}+\ii\sin{x}=y$，那么对应的$\ee^{-\ii x}=\cos{x}-\ii\sin{x}=\frac{1}{y}$

    这样的话，我们就可以用$y$来表示$\cos{x},\ \ii\sin{x}$了
    
    反解一下我们可以得到
    \formula{2\cos{x}=y+\frac{1}{y},\ 2\ii\sin{x}=y-\frac{1}{y}}

    这是针对于$x$，那么如果是$nx$呢

    我们同样操作可得
    \formula{\ee^{\ii nx}=\cos{nx}+\ii\sin{nx}}
    \formula{\ee^{\-ii nx}=\cos{nx}-\ii\sin{nx}}
    而我们知道$\ee^{\ii nx}=\qty(\ee^{ix})^n=y^n,\ \ee^{\-ii nx}=\qty(\ee^{-ix})^n=\frac{1}{y^n}$，因此
    \formula{\cos{nx}+\ii\sin{nx}=y^n}
    \formula{\cos{nx}-\ii\sin{nx}=\frac{1}{y^n}}

    同样的反解出$2\cos{nx}=y^n+\frac{1}{y},\ 2\ii\sin{nx}=y^n-\frac{1}{y^n}$

    到这，是不是发现，我们构建了一个桥梁，高次到低次的桥梁，虽然他没那么明显，不过我们通过具体实操$\int{\sin^8{x}\d x}$这个积分，就知道了.

    首先我们发现，我们刚刚得到的是$2\ii\sin{x}$，因此我们要先配凑一下
    \formula{\int{\sin^8{x}\d x}=\int{\frac{\qty(2\ii\sin{x})^8}{2^8\ii^8}\d x}}

    根据刚刚的公式我们可以得到
    \formula{\qty(2\ii\sin{x})^8=\qty(y-\frac{1}{y})^8}

    通过二项式定理我们可以展开得到
    \formula{
        \begin{aligned}
            \qty(2\ii\sin{x})^8&=\qty(y-\frac{1}{y})^8\\
            &=y^8-C_8^1y^6+C_8^2y^4-C_8^3y^2+C_8^4\\
            &+\frac{1}{y^8}-C_8^1\frac{1}{y^6}+C_8^2\frac{1}{y^4}-C_8^3\frac{1}{y^2}\\
            &=\qty(y^8+\frac{1}{y^8})-8\qty(y^6+\frac{1}{y^6})+28\qty(y^4+\frac{1}{y^4})-56\qty(y^2+\frac{1}{y^2})+70
        \end{aligned}
    }
    
    虽然$y$的结构非常好看，但$y$终究是工具人和桥梁，不可过多留恋.（悲伤）（什么，你不想做工具人，那就速速努力学习，提升价值！）偏了偏了，我们继续回来.

    那么我们再用$\cos{x}=y+\frac{1}{y},\ \cos{nx}=y^n+\frac{1}{y^n}$换回来，可得
    \formula{
        \begin{aligned}
        \qty(2\ii\sin{x})^8&=\qty(y^8+\frac{1}{y^8})-8\qty(y^6+\frac{1}{y^6})+28\qty(y^4+\frac{1}{y^4})-56\qty(y^2+\frac{1}{y^2})+70\\
        &=2\cos{8x}-16\cos{6x}+56\cos{4x}-112\cos{2x}+70
        \end{aligned}
        }

        而$\qty(2\ii\sin{x})^8=2^8\sin^8{x}$，因此我们将高次变成了低次.之后的积分就方便了.
        \formula{
            \begin{aligned}
                \int{\sin^8{x}\d x}&=\int{\frac{\qty(2\ii\sin{x})^8}{2^8\ii^8}\d x}\\
                &=\frac{1}{2^8\ii^8}\int{2\cos{8x}-16\cos{6x}+56\cos{4x}-112\cos{2x}+70}\\
                &=\frac{1}{2^7}\int{\qty(\cos{8x}-8\cos{6x}+28\cos{4x}-56\cos{2x}+35)\d x}\\
                &=\frac{1}{2^7}\qty(\frac{\sin{8x}}{8}-8\frac{\sin{6x}}{6}+28\frac{\sin{4x}}{4}-56\frac{\sin{2x}}{2}+35x)+C
            \end{aligned}
        }

        优雅!实在是太优雅了!倍角法相比于逐级降次的方法来说，计算量和复杂度降低了太多，太优雅了!不过细心的盲点可能发现了华生：“诶，那如果是奇数次的话$\ii$是不是就消不掉了.”说得好！确实对于$\sin^m{x}$的问题，只适用于$m$为偶数的情况，但如果是奇数其实也是可以处理的，后面我们再做探索.而对于$\cos^n{x}$，就不存在奇偶关系了，大家可以自己试试$\sim$今天就先到这了.
}%在花括号内输入分析